\section{Background: Current Practice and Existing Guidance}
\label{sec:background}

\subsection{What Generative Models Do in HRI}

Within a few years, pretrained generative models have been put to almost
every use that a robot architecture or an HRI study has for computation.
They generate a social robot's speech, as in early open-dialogue systems for
Pepper, NAO, and Furhat \cite{billing2023language,cherakara2023furchat} and
in companion robots studied with older adults
\cite{irfan2024recommendations,irfan2025between,lima2025promoting}. They
generate expressive behaviour: gesture and motion as code
\cite{mahadevan2024generative}, facial expressions
\cite{antony2025xpress,mishra2023realtime}, and the body language of legged
robots \cite{roy2025gpt}. They decide: which skill a robot should execute
\cite{ahn2022do,liang2023code}, whether to help or stay silent in a group
\cite{tanneberg2024help}, whether and how to initiate an interaction
\cite{sasabuchi2025agreeing}, when to take a conversational turn
\cite{skantze2025applying}, and whether a path through a crowd is socially
acceptable \cite{song2025vlm}. They perceive: vision-language models give a
shop robot awareness of its customers \cite{edirisinghe2026welcome} and read
nonverbal cues of readiness to be served \cite{bu2025using}. They let end
users program robots in natural language \cite{karli2024alchemist}. In
vision-language-action models they act directly, mapping images and
instructions to motor commands
\cite{zitkovich2023rt2,kim2025openvla,black2025pi0}; surveys of foundation
models in robotics describe the breadth of these uses
\cite{firoozi2025foundation,kawaharazuka2024real}.

Models have entered the research pipeline as well. Language and multimodal
models have been used to annotate users' enjoyment of conversations with a
robot \cite{pereira2024multimodal}, to conduct thematic analysis of
interviews from an HRI study with disabled participants
\cite{markelius2026evaluating}, to cluster and label what people disclose to
robots \cite{chiang2025selfdisclosure}, to stand in for the participants of
published HRI studies \cite{wachowiak2024large}, to complete HRI
questionnaires on a robot's behalf \cite{lockwood2026when}, and to serve as
models of human preference inside planners \cite{zhang2023largeb}.

\subsection{How They Are Currently Evaluated}

A systematic review of 86 LLM-driven HRI studies gives the clearest picture
of current practice \cite{wang2026how}. Questionnaires were used in 70 of
the 86 papers and metrics of perceptual and relational experience in 65,
whereas metrics specific to the performance of the language model appeared
in 30. The reviewers characterise the field as largely exploratory, with
setups, methods, and metrics that vary widely from study to study. A second
review, of 29 studies of LLM-powered robots, likewise finds that measured
outcomes are mostly proximal and recommends that studies report the
consistency of the robot's behaviour, the scope and provenance of its
memory, and its policy for responding and refusing
\cite{nikolova2026systematic}. In the terms we introduce in
Section~\ref{sec:framework}, evidence in HRI is gathered predominantly at
the level of the interaction, and measurement aimed at the model itself is
reported in about a third of studies.

Individual studies span a wide range. At one end are system demonstrations
that report no evaluation of the model's outputs, which is appropriate to
their purpose \cite{billing2023language,cherakara2023furchat}. In the
middle are user studies in which the model is apparatus and the evidence
consists of questionnaires and interviews
\cite{kim2024understanding,rosen2024previous}. At the other end are studies
that show what fuller evaluation looks like, and they are worth naming
because they demonstrate that the practices proposed in this article are
feasible within ordinary HRI research:

\begin{itemize}[leftmargin=1.4em]
\item \emph{Configuration reporting.} Some studies report a dated model
snapshot, the temperature, and a fixed seed ``to ensure reproducibility''
\cite{tanneberg2024help}, a full table of decoding parameters
\cite{mishra2023realtime}, or complete prompts with the observation that
outputs vary even at temperature zero \cite{mahadevan2024generative}. Others
name only a product family.
\item \emph{Coded outputs with agreement statistics.} In a study of a
companion robot with older adults, two coders analysed the conversation
transcripts and reported their agreement, which was moderate for topics and
low for dialogue disruptions \cite{irfan2025between}. A study of a robot in
a geriatric day hospital had two researchers code video independently and
stated that no agreement statistic was computed \cite{blavette2025integrating}.
Both are more informative than silence.
\item \emph{Pipeline measurement.} A study of a conversational robot
reported response latency, the language model's share of that latency, the
deletion rate of the speech recogniser, and technical failures per turn,
with confidence intervals, alongside usability and perception scales
\cite{satake2026focus}. In a care home, speech recognition was 95\% accurate
for the standard language and 45 to 56\% for regional dialects, and the
language model restored coherence after roughly half of the misrecognitions
\cite{pinto2025speech}; with second-language learners, the language model
sustained task conversations despite poor transcripts
\cite{verhelst2025large}. Fluent replies, in other words, can mask an
upstream failure.
\item \emph{Stress and perturbation testing.} Apparent theory-of-mind
reasoning about robot behaviour was shown to collapse under perturbation
tests, with ten repeated runs at each of three temperatures
\cite{verma2024theory}; language models were audited for discrimination and
for acceptance of harmful instructions in robot-task framings
\cite{hundt2025llm}; and students were invited to manipulate an LLM-driven
robot into violating stated principles \cite{abbo2025can}.
\item \emph{Benchmarks built with stakeholders.} Caregiving scenarios
validated by clinical stakeholders \cite{liu2026careeval}, privacy scenarios
for home robots answered by crowd participants and by ten language models,
whose agreement with people was generally low
\cite{sullivan2025benchmarking}, and a benchmark of human perception for
vision-language models that treats latency as a metric
\cite{shi2025hribench}.
\end{itemize}

The same literature documents why the model's outputs need to be examined
directly. Commercially available care software running on a humanoid robot
falsely claimed that it could remind patients to take medication
\cite{ranisch2025rapid}. Primary-school pupils interacting with LLM-powered
tutor and peer robots that had been programmed to make noticeable mistakes
often accepted the wrong answers \cite{rutatola2026robot}. Language models
completing HRI questionnaires were consistent and, on one dimension,
systematically inverted relative to people \cite{lockwood2026when}. And a
model-generated thematic analysis of interviews with disabled participants
agreed with the human analysis only modestly above chance, with codes that
were semantically close to the human analysts' codes nonetheless replacing
participants' own terms, such as \emph{ableism}, with more generic ones
\cite{markelius2026evaluating}.

\subsection{Existing Guidance}

HRI has a substantial methodological literature on which this article
builds. It includes guidance on designing, analysing, and reporting
experiments \cite{bethel2010review,baxter2016characterising,hoffman2020primer};
reporting guidelines for Wizard-of-Oz studies, derived from a review which
found that only 3.7\% of such studies measured wizard error and 5.4\%
reported wizard training \cite{riek2012wizard}; audits of statistical
practice, which found that 4 of 144 papers applied appropriate tests to
correctly designed rating scales \cite{schrum2023concerning} and that most
papers reported no power analysis \cite{bartlett2022have}; analyses of
replicability \cite{leichtmann2022crisis,gunes2022reproducibility,
stower2026children}; reviews of who takes part in HRI studies and how that
is reported \cite{seaborn2023not,cordero2022review}; and guidance on the
choice of measurement scales \cite{saad2026choosing}. Standard reporting
forms have been proposed \cite{bagchi2023towards,holthaus2025would}, and a
consensus-based reporting checklist for HRI user studies has recently been
produced through a Delphi process with 34 experts \cite{chandra2026rush}.

Guidance specific to generative models is recent and concerns language
models in the robot. Williams and colleagues characterised LLMs in robot
architectures as ``Scarecrows'', black-box modules that quickly enable
full-pipeline systems much as Wizard-of-Oz techniques do, and proposed
initial reporting guidelines modelled on those for Wizard-of-Oz studies
\cite{williams2024scarecrows}. A subsequent article in this journal
identifies five groups of stakeholders in HRI research, what each needs to
know about the use of an LLM, and suggested guidelines for what to report
\cite{matuszek2026reporting}. Prompting guidelines with a structured
template address how a robot's persona and limits are specified
\cite{ashok2026why}; a framework for evaluating LLM-generated language in
robot role-play argues that factuality-style evaluation is insufficient
\cite{hohn2024using}; and measures have been proposed for sycophancy in
social robots \cite{seaborn2026robotic}.

The present article is complementary to this work in three respects. First,
its subject is evaluation and not only reporting: what evidence about a
model should be gathered, and how, before the question of how to report it
arises. Second, it differentiates by role, on the argument that reporting
and evaluation requirements for a model that speaks, a model that embeds,
and a model that moves an arm are different. Third, it extends the scope
from language models in the robot to vision-language and
vision-language-action models, and to models used in the research pipeline
as annotators, judges, and simulators, which existing guidance does not
address.

\subsection{Three Traditions}

The practices proposed below are drawn from three bodies of methodological
work, each of which has confronted generative models on its own terms.

\emph{Natural language processing} has a long experience of evaluating
generated text with human raters and of the weaknesses of that practice
\cite{vanderlee2021human,howcroft2020twenty,gehrmann2023repairing}, a more
recent and fast-growing literature on using models as judges
\cite{zheng2023judging,gu2026survey}, and an evaluation science that has
begun to treat the measurement of generative systems as a problem of
validity \cite{jacobs2021measurement,wallach2025position,bean2025measuring}
and to study interaction and not only output
\cite{lee2023evaluating,ibrahim2025towards}.

\emph{Psychology and content analysis} supply the theory and statistics of
agreement among raters \cite{krippendorff2019content,hallgren2012computing},
the theory of validity \cite{cronbach1955construct,messick1995validity,
kane2013validating}, the reforms that followed the replication crisis
\cite{opensciencecollaboration2015estimating,simmons2011falsepositive,
nosek2018preregistration}, and, in the last three years, a critical
literature on language models as annotators
\cite{gilardi2023chatgpt,pangakis2023automated,baumann2025large} and as
substitutes for human participants \cite{argyle2023out,wang2025large}.

\emph{Robotics} supplies the practice of evaluating learned policies by
physical trials and the recent effort to make that practice statistically
sound \cite{kressgazit2024robot,agarwal2021deep}, simulation and distributed
evaluation for generalist policies \cite{li2025evaluating,atreya2025roboarena},
and the analysis of safety for robots controlled by language models
\cite{robey2025jailbreaking,sermanet2025generating,hundt2025llm}.

None of the three is sufficient for HRI. NLP's methods assume text read on a
screen. Psychology's assume a human rater or a fixed instrument. Robotics'
assume that success is a property of the task and not of the person on its
receiving end. The framework that follows is an attempt to say which
tradition governs which use of a model, and what must be added when the use
is embodied and social.
